% =====================================================================
%  SU(2) e SO(3) — o homomorfismo de recobrimento duplo
%  Notas de aula (nível: início de mestrado em Física Teórica)
%  Transcrição e edição de notas manuscritas.
%
%  CONVENÇÃO DE EDIÇÃO
%  -------------------
%  Texto em PRETO  ........ transcrição fiel das notas originais.
%  Texto em VERMELHO (com o adereço ▸) ... acréscimos, correções e
%                          passagens de rigor introduzidas nesta edição.
%  Para "aceitar" todas as mudanças de uma só vez, basta trocar a cor
%  de destaque para preto no painel abaixo (corDestaque -> {0,0,0}).
% =====================================================================

\documentclass[12pt,a4paper]{article}

% ---------- codificação e idioma ----------
\usepackage[T1]{fontenc}      % fontes EC = Computer Modern na codificação T1
\usepackage[utf8]{inputenc}
\usepackage[brazil]{babel}

% ---------- Computer Modern ----------
% A fonte-padrão do LaTeX JÁ É a Computer Modern.
% (NÃO carregar 'lmodern' nem 'times'/'newtx' para preservá-la.)

% ---------- matemática ----------
\usepackage{amsmath,amssymb,mathtools,amsthm}

% ---------- cor, layout, links ----------
\usepackage{xcolor}
\usepackage{geometry}
\usepackage{setspace}
\usepackage{indentfirst}
\usepackage{enumitem}
\usepackage[hidelinks]{hyperref}

% =====================================================================
%  ██  PAINEL DE FORMATAÇÃO — edite livremente aqui  ██
% =====================================================================

% (1) Margens ABNT: esquerda/superior 3 cm, direita/inferior 2 cm.
\geometry{left=3cm, top=3cm, right=2cm, bottom=2cm}

% (2) Espaçamento entre linhas (ABNT = 1,5).
%     Troque por \singlespacing ou \doublespacing se quiser.
\onehalfspacing

% (3) Recuo de primeira linha do parágrafo (ABNT ~ 1,25 cm).
\setlength{\parindent}{1.25cm}

% (4) Cor dos acréscimos/correções desta edição.
%     >>> Mude para {0,0,0} (preto) para incorporar tudo ao texto. <<<
\definecolor{corDestaque}{RGB}{190,25,25}

% (5) Adereço (marcador) que abre cada trecho editado.
\newcommand{\marcador}{\ensuremath{\color{corDestaque}\blacktriangleright}\,}

% =====================================================================
%  Comandos de destaque das edições
% =====================================================================
% \add{...}      -> trecho curto, em linha, em vermelho, iniciado por ▸
% ambiente 'nota'-> bloco inteiro em vermelho, iniciado por ▸
\newcommand{\add}[1]{{\marcador\color{corDestaque}#1}}
\newenvironment{nota}{\par\color{corDestaque}\noindent\marcador\ignorespaces}{\par}

% =====================================================================
%  Abreviações matemáticas (edite/expanda conforme preferir)
% =====================================================================
\newcommand{\C}{\mathbb{C}}
\newcommand{\R}{\mathbb{R}}
\newcommand{\Z}{\mathbb{Z}}
\newcommand{\id}{\mathbb{1}}                 % matriz identidade
\newcommand{\dg}{\dagger}
\newcommand{\Tr}{\operatorname{Tr}}
\newcommand{\su}{\mathfrak{su}}
\newcommand{\so}{\mathfrak{so}}
\newcommand{\vsig}{\vec{\sigma}}
\newcommand{\vx}{\vec{x}}
\newcommand{\vn}{\hat{n}}
\newcommand{\comm}[2]{\left[#1,#2\right]}
\newcommand{\acomm}[2]{\left\{#1,#2\right\}}

% ---------- ambientes de enunciado (estilo enxuto) ----------
\theoremstyle{definition}
\newtheorem{definicao}{Definição}
\theoremstyle{plain}
\newtheorem{proposicao}{Proposição}
\theoremstyle{remark}
\newtheorem*{observacao}{Observação}
\renewcommand{\qedsymbol}{$\blacksquare$}

% =====================================================================
\title{\textbf{$SU(2)$ e $SO(3)$:\\[2pt] o homomorfismo de recobrimento duplo}}
\author{Notas de aula\thanks{Transcrição editada de notas manuscritas. Nível: início de mestrado em Física Teórica.}}
\date{\today}

\begin{document}
\maketitle

\begin{nota}
\emph{Sobre esta edição.} O texto em preto reproduz as notas originais; o
texto em vermelho, aberto pelo adereço \marcador, assinala acréscimos,
correções e passagens de rigor introduzidas para tornar a leitura
autossuficiente. Para incorporar todas as alterações ao corpo principal,
basta trocar \texttt{corDestaque} para preto no painel de formatação.
\end{nota}

% =====================================================================
\section{Introdução}
% =====================================================================
\begin{nota}
Os grupos $SU(2)$ e $SO(3)$ têm a \emph{mesma} álgebra de Lie e, portanto,
descrevem localmente as mesmas rotações; ainda assim, não são isomorfos. O
objetivo destas notas é construir explicitamente um homomorfismo
sobrejetor $R:SU(2)\to SO(3)$, determinar seu núcleo e concluir o
isomorfismo
\[
  SO(3)\;\cong\;SU(2)/\Z_2 .
\]
Essa relação de recobrimento duplo (\emph{double cover}) é a origem
algébrica do \emph{spin} semi-inteiro e do fato de que um elétron precisa
girar $4\pi$ — e não $2\pi$ — para retornar ao estado inicial.
\end{nota}

% =====================================================================
\section{O grupo $SU(2)$}
% =====================================================================

\begin{definicao}
O grupo unitário especial de grau $2$ é
\[
  SU(2):=\bigl\{\,U\in M_2(\C)\ :\ U^{\dg}U=\id,\ \det U=1\,\bigr\}.
\]
\end{definicao}

Impondo $U^{\dg}U=\id$ a uma matriz genérica, sua inversa é a adjunta.
Comparando
\[
  U^{-1}=\begin{pmatrix} d & -b\\ -c & a\end{pmatrix}
  =\begin{pmatrix} \bar a & \bar c\\ \bar b & \bar d\end{pmatrix}
  \;\Longrightarrow\;
  U=\begin{pmatrix} a & b\\ -\bar b & \bar a\end{pmatrix},
\]
de modo que todo elemento se escreve como
\begin{equation}\label{eq:formaU}
  SU(2)=\left\{
    \begin{pmatrix} a & b\\ -\bar b & \bar a\end{pmatrix}
    \ :\ a,b\in\C,\ |a|^2+|b|^2=1
  \right\},
\end{equation}
pois $\det U=|a|^2+|b|^2=1$.

Separando parte real e imaginária, $a=a_1+ia_2$ e $b=b_1+ib_2$, obtemos
\begin{equation}\label{eq:Upauli}
  U=\begin{pmatrix} a_1+ia_2 & b_1+ib_2\\ -b_1+ib_2 & a_1-ia_2\end{pmatrix}
   = a_1\,\id + i\,\vec b\cdot\vsig ,
\end{equation}
onde $\vec b=(b_2,b_1,a_2)$ e $\vsig=(\sigma_1,\sigma_2,\sigma_3)$ são as
matrizes de Pauli
\[
  \sigma_1=\begin{pmatrix}0&1\\1&0\end{pmatrix},\quad
  \sigma_2=\begin{pmatrix}0&-i\\ i&0\end{pmatrix},\quad
  \sigma_3=\begin{pmatrix}1&0\\0&-1\end{pmatrix}.
\]

\begin{observacao}
A condição $a_1^2+a_2^2+b_1^2+b_2^2=1$ identifica $SU(2)$, como variedade,
com a esfera $S^3\subset\R^4$. \add{Em particular, $SU(2)$ é conexo e
simplesmente conexo — fato que usaremos adiante.}
\end{observacao}

As matrizes de Pauli satisfazem a relação fundamental
\begin{equation}\label{eq:pauliprod}
  \sigma_i\sigma_j=\delta_{ij}\,\id+i\,\varepsilon_{ijk}\,\sigma_k,
\end{equation}
da qual seguem, tomando parte antissimétrica e simétrica,
\begin{align}
  \comm{\sigma_i}{\sigma_j}&=2i\,\varepsilon_{ijk}\,\sigma_k,\label{eq:commpauli}\\
  \acomm{\sigma_i}{\sigma_j}&=2\,\delta_{ij}\,\id.\label{eq:acommpauli}
\end{align}
\add{Delas decorrem ainda $\Tr\sigma_i=0$ e a relação de ortogonalidade}
\begin{equation}\label{eq:traco}
  \add{\Tr(\sigma_i\sigma_j)=2\,\delta_{ij},}
\end{equation}
\add{que será a ferramenta para extrair componentes de uma matriz.}

% =====================================================================
\section{A correspondência vetor--matriz}
% =====================================================================

A cada vetor $\vx=(x_1,x_2,x_3)\in\R^3$ associamos a matriz
\begin{equation}\label{eq:X}
  X:=\vx\cdot\vsig=
  \begin{pmatrix} x_3 & x_1-ix_2\\ x_1+ix_2 & -x_3\end{pmatrix}.
\end{equation}
Ela é \emph{hermitiana} ($X^{\dg}=X$) e tem \emph{traço nulo}
($\Tr X=0$). Além disso,
\begin{equation}\label{eq:detX}
  \det X=-x_3^2-(x_1^2+x_2^2)=-\,|\vx|^2 .
\end{equation}

\begin{nota}
A correspondência $\vx\mapsto X$ é um isomorfismo linear entre $\R^3$ e o
espaço vetorial real das matrizes $2\times2$ hermitianas de traço nulo, que
tem dimensão $3$ e admite $\{\sigma_1,\sigma_2,\sigma_3\}$ como base. A
inversa é dada explicitamente por
\[
  x_i=\tfrac12\,\Tr(\sigma_i X),
\]
como se verifica usando \eqref{eq:traco}. Esse ponto — omitido nas notas —
é o que garante que qualquer matriz hermitiana de traço nulo possa ser
reescrita na forma $\vx\cdot\vsig$.
\end{nota}

% =====================================================================
\section{A ação de $SU(2)$ e a rotação induzida}
% =====================================================================

Dado $U\in SU(2)$, definimos a ação por conjugação
\begin{equation}\label{eq:acao}
  X\;\longmapsto\;X'=U\,X\,U^{\dg}.
\end{equation}

\begin{proposicao}\label{prop:invariantes}
Se $X$ é hermitiana de traço nulo, então $X'=UXU^{\dg}$ também o é, e
$\det X'=\det X$.
\end{proposicao}
\begin{proof}
Hermiticidade e traço:
\begin{gather*}
  (X')^{\dg}=(UXU^{\dg})^{\dg}=U X^{\dg} U^{\dg}=UXU^{\dg}=X',\\[2pt]
  \Tr X'=\Tr(UXU^{\dg})=\Tr(X\,U^{\dg}U)=\Tr X=0 .
\end{gather*}
Determinante, usando $\det U=1$:
\[
  \det X'=\det U\,\det X\,\det U^{\dg}=\det X. \qedhere
\]
\end{proof}

Pela Proposição~\ref{prop:invariantes}, $X'$ é novamente da forma
$X'=\vx'\cdot\vsig$ para algum $\vx'\in\R^3$, e de \eqref{eq:detX} com
$\det X'=\det X$ segue
\begin{equation}\label{eq:norma}
  |\vx'|^2=|\vx|^2 .
\end{equation}

\begin{nota}
A aplicação $\vx\mapsto\vx'$ é \emph{linear}: a ação \eqref{eq:acao} é linear
em $X$ e as passagens $\vx\leftrightarrow X$ o são. Uma transformação linear
de $\R^3$ que preserva a norma \eqref{eq:norma} é ortogonal; logo existe
$R(U)\in O(3)$ com
\[
  \vx'=R(U)\,\vx .
\]
Ainda falta mostrar que $\det R(U)=+1$, i.e.\ que $R(U)\in SO(3)$
(Seção~\ref{sec:hom}).
\end{nota}

Combinando $x'_i=\tfrac12\Tr(\sigma_i X')$ com $X'=U(x_j\sigma_j)U^{\dg}$,
obtemos a forma explícita das componentes de $R(U)$:
\begin{equation}\label{eq:Rij}
  \boxed{\,R_{ij}(U)=\tfrac12\,\Tr\!\bigl(\sigma_i\,U\,\sigma_j\,U^{\dg}\bigr).\,}
\end{equation}

% =====================================================================
\section{$R$ é homomorfismo e $\det R=1$}\label{sec:hom}
% =====================================================================

\begin{proposicao}
A aplicação $R:SU(2)\to O(3)$ definida por \eqref{eq:Rij} é um
homomorfismo de grupos.
\end{proposicao}
\begin{proof}
Para $U_1,U_2\in SU(2)$,
\[
  (U_1U_2)\,X\,(U_1U_2)^{\dg}=U_1\bigl(U_2 X U_2^{\dg}\bigr)U_1^{\dg},
\]
isto é, conjugar por $U_1U_2$ equivale a conjugar por $U_2$ e depois por
$U_1$. Traduzindo para $\R^3$, $R(U_1U_2)=R(U_1)R(U_2)$.
\end{proof}

\begin{proposicao}
$\det R(U)=1$ para todo $U\in SU(2)$; portanto $R(U)\in SO(3)$.
\end{proposicao}
\begin{proof}
\add{A aplicação $U\mapsto\det R(U)$ é contínua e toma valores no conjunto
discreto $\{+1,-1\}$ (pois $R(U)\in O(3)$). Como $SU(2)$ é conexo, sua
imagem por uma função contínua é conexa; logo $\det R(U)$ é constante. Em
$U=\id$ temos $R(\id)=\id$ e $\det R(\id)=1$, donde $\det R(U)=1$ para
todo $U$.}
\end{proof}

% =====================================================================
\section{Núcleo, sobrejetividade e o isomorfismo}
% =====================================================================

\subsection*{Núcleo}
Por definição,
\[
  \ker R=\{\,U\in SU(2)\ :\ R(U)=\id\,\}
        =\{\,U\ :\ UXU^{\dg}=X\ \ \forall\,X\,\},
\]
ou seja, $U$ comuta com $\sigma_1,\sigma_2,\sigma_3$.
\add{Como as matrizes de Pauli, junto com $\id$, geram toda a álgebra
$M_2(\C)$, comutar com todas elas obriga $U$ a comutar com \emph{qualquer}
matriz $2\times2$; pelo Lema de Schur, $U=\lambda\id$.} De $\det U=1$ segue
$\lambda^2=1$, isto é, $\lambda=\pm1$. Logo
\begin{equation}\label{eq:kernel}
  \ker R=\{+\id,\,-\id\}\cong\Z_2 .
\end{equation}
Em particular $R(U)=R(-U)$: a aplicação é \emph{dois-para-um}.

\subsection*{Sobrejetividade}
\begin{proposicao}\label{prop:exp}
Para $\vn\in\R^3$ unitário e $\theta\in\R$,
\begin{equation}\label{eq:expU}
  U(\theta,\vn)=\exp\!\Bigl(-\tfrac{i\theta}{2}\,\vn\cdot\vsig\Bigr)
  =\cos\tfrac{\theta}{2}\,\id-i\,\sin\tfrac{\theta}{2}\,\vn\cdot\vsig
  \ \in SU(2),
\end{equation}
e $R\bigl(U(\theta,\vn)\bigr)$ é a rotação de ângulo $\theta$ em torno do
eixo $\vn$.
\end{proposicao}
\begin{nota}
A expansão em \eqref{eq:expU} usa $(\vn\cdot\vsig)^2=\id$, consequência de
\eqref{eq:pauliprod} com $|\vn|=1$. Como $\theta$ e $\vn$ são arbitrários e
toda rotação de $SO(3)$ é uma rotação por um ângulo em torno de um eixo
(teorema de Euler), a Proposição~\ref{prop:exp} mostra que $R$ é
sobrejetora. A identificação explícita de $R(U)$ com uma rotação é feita no
exemplo da Seção~\ref{sec:exemplo}.
\end{nota}

\subsection*{O isomorfismo}
Com $R$ homomorfismo sobrejetor e $\ker R\cong\Z_2$, o teorema fundamental
do isomorfismo fornece
\begin{equation}\label{eq:iso}
  \boxed{\,SU(2)/\Z_2\;\cong\;SO(3).\,}
\end{equation}

% =====================================================================
\section{Exemplo: rotação em torno de $z$}\label{sec:exemplo}
% =====================================================================

Tomando $\vn=\hat z=(0,0,1)$ em \eqref{eq:expU},
\[
  U=\cos\tfrac{\theta}{2}\,\id-i\sin\tfrac{\theta}{2}\,\sigma_3
   =\begin{pmatrix} e^{-i\theta/2} & 0\\ 0 & e^{i\theta/2}\end{pmatrix}.
\]
Conjugando $X=\vx\cdot\vsig$,
\[
  X'=UXU^{\dg}
    =\begin{pmatrix} x_3 & e^{-i\theta}(x_1-ix_2)\\
                     e^{i\theta}(x_1+ix_2) & -x_3\end{pmatrix},
\]
de onde se lê, comparando com \eqref{eq:X},
\[
  \begin{aligned}
    x_1'&=x_1\cos\theta-x_2\sin\theta,\\
    x_2'&=x_1\sin\theta+x_2\cos\theta,\\
    x_3'&=x_3,
  \end{aligned}
  \qquad\Longleftrightarrow\qquad
  R(U)=\begin{pmatrix}
    \cos\theta & -\sin\theta & 0\\
    \sin\theta & \phantom{-}\cos\theta & 0\\
    0 & 0 & 1
  \end{pmatrix},
\]
exatamente a rotação de ângulo $\theta$ em torno de $\hat z$.

% =====================================================================
\section{Recobrimento duplo e álgebra de Lie}
% =====================================================================

\subsection*{Recobrimento duplo}
Substituindo $\theta\to\theta+2\pi$ em \eqref{eq:expU},
\[
  U(\theta+2\pi,\vn)
  =\cos\!\Bigl(\tfrac{\theta}{2}+\pi\Bigr)\id
   -i\sin\!\Bigl(\tfrac{\theta}{2}+\pi\Bigr)\vn\cdot\vsig
  =-\,U(\theta,\vn),
\]
enquanto a rotação não muda: $R\bigl(U(\theta+2\pi)\bigr)=R(-U)=R(U)$.
É preciso variar $\theta$ de $0$ a $4\pi$ para que $U$ retorne ao valor
inicial, ao passo que $R$ já se repete em $2\pi$. Esse é o sentido preciso
de $SU(2)$ ser o \emph{recobrimento duplo} de $SO(3)$.

\begin{nota}
Topologicamente, $SU(2)\cong S^3$ é simplesmente conexo, enquanto
$SO(3)\cong \R P^3$ tem grupo fundamental $\pi_1(SO(3))\cong\Z_2$. A
projeção $S^3\to\R P^3$ que identifica pontos antípodas ($U\sim -U$) é
precisamente a aplicação $R$.
\end{nota}

\subsection*{Isomorfismo das álgebras de Lie}
Os geradores anti-hermitianos de $\su(2)$,
\[
  J_a=-\,\frac{i}{2}\,\sigma_a,
  \qquad
  \comm{J_a}{J_b}=\varepsilon_{abc}\,J_c,
\]
têm as \emph{mesmas} constantes de estrutura que os geradores de $\so(3)$,
$(T_a)_{bc}=-\varepsilon_{abc}$, os quais satisfazem
$\comm{T_a}{T_b}=\varepsilon_{abc}T_c$. A aplicação linear
$\varphi:\su(2)\to\so(3)$, $\varphi(J_a)=T_a$, preserva o colchete,
\[
  \varphi\bigl(\comm{J_a}{J_b}\bigr)=\comm{\varphi(J_a)}{\varphi(J_b)},
\]
e é bijetora; portanto
\begin{equation}\label{eq:isoalg}
  \su(2)\cong\so(3).
\end{equation}
\add{As álgebras são isomorfas (os grupos são localmente idênticos), mas os
\emph{grupos} não o são: a diferença é inteiramente global, capturada por
$\ker R\cong\Z_2$ em \eqref{eq:kernel}.}

% =====================================================================
\section{Consequências físicas}
% =====================================================================
\begin{itemize}[leftmargin=1.6em,itemsep=2pt]
  \item As representações irredutíveis de $SU(2)$ são rotuladas por
        $j=0,\tfrac12,1,\tfrac32,\dots$. As de spin semi-inteiro
        \emph{não} são representações (uni-valentes) de $SO(3)$: sob elas,
        $U$ e $-U$ agem distintamente, embora correspondam à mesma rotação.
  \item Um estado de spin $\tfrac12$ adquire fase $-1$ sob rotação de
        $2\pi$ e só retorna a si mesmo após $4\pi$ — manifestação direta de
        $U(2\pi,\vn)=-\id$.
\end{itemize}

% =====================================================================
\begin{nota}
\section*{Referências}
\addcontentsline{toc}{section}{Referências}
As referências abaixo, todas padrão para o tema, foram acrescentadas nesta
edição.
\begin{itemize}[leftmargin=1.6em,itemsep=2pt]
  \item HALL, B. C. \emph{Lie Groups, Lie Algebras, and Representations}.
        2.~ed. Springer, 2015.
  \item GEORGI, H. \emph{Lie Algebras in Particle Physics}.
        2.~ed. Westview Press, 1999.
  \item SAKURAI, J. J.; NAPOLITANO, J. \emph{Modern Quantum Mechanics}.
        2.~ed. Cambridge University Press, 2017.
  \item TUNG, W.-K. \emph{Group Theory in Physics}. World Scientific, 1985.
\end{itemize}
\end{nota}

\end{document}
